\section{模型假设}

1、假设同一患者编号与同一精确开单时刻对应一次完整检查申请；医嘱字段中以“][”连接的内容表示分别计费的项目，应拆为同一事件下的多个任务。该假设直接对应源系统字段结构。

2、假设题目给出的项目时长区间能够代表标准检查工时；名义场景取区间中点，题目所述5\%复杂检查取区间上界。由于附件没有复杂病例标签，主结果用固定随机种子抽取，并用5个额外种子检验。

3、假设工作模板为每日08:00--12:00和13:00--17:00，检查不得跨越午休或日界。该模板由训练期医生首末活动记录重构，是标准日班压力场景，不解释为医院正式排班表。

4、假设任一项目同时占用一台兼容设备和一名医生；设备独占，医生以星期—5分钟槽总并发容量表示。题面33名未来医生与历史ID无法一一映射，因此不虚构实名医生专项资质矩阵。

5、假设项目—设备兼容只由题目表1的项目文本确定；允许直接名称、明确同义或包含关系以及题面明确的广义项目描述，复合项目取组成项目的同设备交集。无法从表1确认的项目不进入候选设备集合。

6、假设床旁表示检查地点而非项目能力。床旁任务只有在项目本身与7号机能力相容时才可使用7号机，床旁标记不能覆盖心脏、血管、产科、介入等专项要求。

7、假设腹部、胃肠和肠系膜等空腹项目须在10:00前完成；憋尿项目在开单后至少准备60分钟。项目准备类型对应公开超声检查流程中的空腹与充盈膀胱要求\cite{hrhospital2025}，具体60分钟提前量由敏感性界定。准备约束按项目粒度施加，事件级统一准备作为对照情景。

8、假设同一患者的两个项目不能重叠；跨设备相邻项目至少留10分钟转运，同设备连续项目不额外计转运。5分钟和15分钟边界另行检验。

9、假设住院事件在开单后48小时内全部项目完成才计为达标；未排入、项目缺失或任一项目越过截止均记为未完成，分母不因能力缺失而减少。

10、假设门诊、体检报告日期可作为回顾性计划日压力代理，但不能还原真实预约日期；其最后项目在计划日17:00前结束记为计划日完成。本文据此讨论服务水平，不声称复原历史调度。

11、假设最后一个排程项目的结束时刻是报告完成的乐观代理。附件没有扫查结束、书写和提交时刻，故模型结果用于预约与容量评估，不能解释为真实报告时效的因果改善。

12、假设训练期结束于2024年3月31日，留出期从2024年4月1日至2025年3月31日。阈值、项目目录、医生容量和稀缺权重只读训练期；留出期仅用于评价，不反向调参。
